\scanpage{14}
Si $E=A$ est une Algèbre qu. coh.\inserted{(\unsure{resp. et graduée})}, il
en est de même de $p_S(A)$, et on aura
\[
\left\{\begin{aligned}
p_S(\operatorname{Spec}(A))&=\operatorname{Spec}(p_S(A))\\
p_S(\operatorname{Proj}(A))&=\operatorname{Proj}(p_S(A))
\end{aligned}\right.
\]
\margin{\diagram{\begin{tikzcd}[column sep=small,row sep=small]
& p_S(X') \arrow[r] \arrow[d] & X' \arrow[d]\\
X \arrow[r] \arrow[dr,"f"'] & p_S(X) \arrow[r] \arrow[d] & X \arrow[d,"f"]\\
& S \arrow[r,"{\operatorname{Frob}(S)}"'] & S
\end{tikzcd}}}
et les hom $\operatorname{Frob}_S(X)$ pour $X=\operatorname{Spec}(A)$,
resp. $X=\operatorname{Proj}(A)$ sont ceux déduits de
$\operatorname{Frob}_S(A):p_S(A)\longrightarrow A$, grâce à
la puissance \unsure{$p$ième}. (\unsure{$q$ième}).

Corollaires :

\deleted{Exemples : $E=\underline{\mathcal O}$}
\[
\left\{\begin{aligned}
p_S(\underset{\sim}{\mathrm V}(\underline E))&=\underset{\sim}{\mathrm V}(p_S(\underline E))\\
p_S(\underset{\sim}{\mathrm P}(\underline E))&=\underline{\mathrm P}(p_S(\underline E))
\end{aligned}\right.
\]
les hom de Frobenius pour les précédents sont ceux associés à
\deleted{$p_S(E)\quad\operatorname{Frob}_S(\underline E):p_S(\underline E)\longrightarrow\underline E$}.
\deleted{\unsure{grâce à hom}}
\[
\underline E^{(p)}\longrightarrow\operatorname{Symm}^p(\underline E)\subset\operatorname{Symm}(\underline E)
\]
déduit de l'élévation à la puissance \unsure{$p$ième}.
\margin{\diagram{\begin{tikzcd}[column sep=small,row sep=large]
X' \arrow[r,"{\operatorname{Frob}_X(X')}"] \arrow[rr,bend left=30,"{\operatorname{Frob}_S(X')}"] \arrow[dr,"g"'] & p_X(X') \arrow[r,"\unsure{r}"] \arrow[d] & p_S(X') \arrow[r] \arrow[d] & X' \arrow[d,"g"]\\
& X \arrow[r,"{\operatorname{Frob}_S(X)}"] \arrow[rr,bend right=18,dashed,"{\operatorname{Frob}(X)}"'] \arrow[dr,"f"'] & p_S\deleted{\unsure{f}}(X) \arrow[r] \arrow[d] & X \arrow[d,"f"]\\
&&S \arrow[r,"{\operatorname{Frob}(S)}"'] & S
\end{tikzcd}}
trois carrés cartésiens}
\margin{\figure{\begin{tikzpicture}
\draw (0,-2.5) .. controls (-.3,-1.4) and (-.05,-.7) .. (-.45,-.5)
.. controls (-.1,-.35) and (-.45,.3) .. (0,1);
\draw[->] (-.45,-.5) .. controls (-.5,1) and (-.5,6) .. (-.45,7)
.. controls (-.45,7.6) and (.3,7.8) .. (.45,8.1);
\end{tikzpicture}}}
[N.B. la structure vectorielle de $\mathrm V(\underline E)$ n'est
pas conservée, mais bien entendu la structure en
groupe est conservée. Noter que dans l'hom de
Frobenius $\underset{\sim}{\mathrm P}(\underline E)\longrightarrow
p_S(\underset{\sim}{\mathrm P}(\underline E))=\underset{\sim}{\mathrm P}(p_S(\underline E))$, corres\unsure{.}
\mspage{13}
